% !TEX root = ../../tree_notes.tex
% =====================================================================
%  Chapter 4 opener -- "Chapter at a Glance", reading paths + ladder
%  staircase, BOTH running examples (binary R reused from Ch.3; general
%  contest tree G introduced here), roadmap, topic table, recognition
%  guide, constraint table, notation.
%  NO numbered boxes here (they would be numbered 4.0.x).
% =====================================================================

\begin{chaptertheorybox}{Chapter at a Glance}
Chapters~2 and~3 treated a tree as a \emph{logical} object -- a root, some
children, an order in which to visit them. But the compiler does not know what a
"child" is. Somewhere in RAM a tree is nothing but \term{numbers and pointers}: a
\code{TreeNode} holding two addresses, or a \code{vector<int> adj[N]} of neighbour
lists, or a flat array where node $i$'s children live at $2i+1$ and
$2i+2$. This chapter is about that \emph{physical} layer -- \term{representation} --
and the three decisions it forces on you:

\begin{itemize}[nosep]
  \item \textbf{How do I hold the edges?} Pointers suit a \emph{fixed-arity} binary
        tree; an \term{adjacency list} suits an unbounded \emph{general} tree; a
        \term{parent array} and an \term{implicit array index} suit complete, dense,
        or heap-shaped data. Each choice changes your constant factors, your cache
        behaviour, and whether a chain is $O(1)$ or $O(n)$ to walk up.  \hfill \lv{1}--\lv{3}
  \item \textbf{Which node is the root?} Contest input gives \emph{undirected edges}.
        A tree only becomes a tree you can recurse on once you \term{root} it -- pick
        a start and orient every edge away from it with a parent-guarded DFS. Change
        the root and every $\parent$, $\depth$ and $\subtree$ changes.   \hfill \lv{3}--\lv{4}
  \item \textbf{How did the judge \emph{serialise} it?} Edge lists, parent sequences,
        bracket strings, and LeetCode's level-order-with-nulls array are
        four encodings of the \emph{same} tree; reading them correctly -- and
        \emph{decoding} one back into pointers -- is a whole family of problems. \hfill \lv{2}--\lv{5}
\end{itemize}

\textbf{You will learn to}
\begin{itemize}[nosep]
  \item choose pointer, adjacency-list, parent-array and implicit-index storage by
        the shape and operation mix a problem implies;                    \hfill\lv{1}--\lv{3}
  \item root an undirected edge list with a parent-guarded DFS/BFS and read off
        $\parent$, $\depth$, $\size$ and $\height$ for every node;         \hfill\lv{3}--\lv{4}
  \item recognise a cheap \emph{root change} (rerooting) when a problem asks the same
        question from every node, and know its $O(n)$ form;               \hfill\lv{4}--\lv{5}
  \item decode the four common contest encodings (edge list, parent array, bracket
        string, level-order array) into a workable structure;              \hfill\lv{2}--\lv{4}
  \item rebuild a tree from a self-delimiting string (a preorder with null marks).
        \hfill\lv{4}--\lv{5}
\end{itemize}

\textbf{Prerequisites.} The adjacency-build loop of Chapter~1 (\coderef{c01-build})
and the \code{TreeNode} model of Chapter~2 (\secref{c02-model}); the four DFS \emph{orders}
of Chapter~3 -- this chapter \emph{uses} those orders to fill arrays, it does not
re-derive them (\chref{c03}).

\textbf{In contests and interviews.} Almost every tree problem on LeetCode hands you
a pointer \code{TreeNode}; almost every tree problem on Codeforces hands you $n{-}1$
undirected edges and a $1$-indexed adjacency list. Being fluent at \emph{both}, and at
converting between them, is the quiet skill that separates a clean solution from a
buggy one. This chapter builds that fluency end to end.
\end{chaptertheorybox}

% ---------------- Station 1b: reading paths + ladder position ----------------
\begin{chaptertheorybox}{How to Read This Chapter (two paths up the ladder)}
Level badges $\lv1$--$\lv6$ as before: $\lv1$ is \emph{``store and walk it''}, $\lv3$
is the solid interview medium, $\lv5$ is LC Hard with a decode or a reroot on top.

\textbf{Beginner path ($\lv1$--$\lv3$).} Read \secref{c04-reps} fully -- it fixes the
four storage models and the running examples' two views. Then
\secref{c04-rooting} up through the parent-guarded DFS and subtree sizes. Take every
$\lv1$--$\lv3$ solved problem. You may \emph{skip on the first pass}: the reroot preview
of \secref{c04-rooting}, the bracket/level-order decoders of \secref{c04-formats}, every
$\lv4{+}$ variation, the hardening notes and the Summit. At the end of this path you
can build any tree from any reasonable input and compute its rooted properties.

\textbf{Full path.} Everything, in order -- it ends with the scaled summit (LC Hard 834
Sum of Distances in Tree: root it, get its subtree sizes, then \emph{reroot} in one
$O(n)$ sweep) and the reduction gallery showing that rooting, subtree aggregation and
the $O(1)$ root-change are the same DFS seen from three sides. The genuine $\lv6$
storage problems -- persistent orders, heavy-light layouts, implicit treaps -- live in
the query and balanced chapters (\chref{c24}, \chref{c06}).
\end{chaptertheorybox}

\begin{figure}[H]
\centering
\begin{tikzpicture}[diagram, x=9mm, y=7mm]
  \foreach \st/\lab in {1/lv1,2/lv2,3/lv3,4/lv4,5/lv5,6/lv6} {
    \fill[\lab!22, draw=\lab, line width=0.8pt] (\st-1,0) rectangle (\st,\st*0.52);
    \node[font=\scriptsize\bfseries, text=\lab] at (\st-0.5,\st*0.52-0.42) {L\st};}
  \node[annot, anchor=south west, font=\scriptsize]
    at (0.05,0.62) {4.1 storage models};
  \node[annot, anchor=south west, font=\scriptsize]
    at (1.05,1.14) {4.1 parent arrays / N-ary};
  \node[annot, anchor=south west, font=\scriptsize]
    at (2.05,1.66) {4.2 rooting + subtree sizes};
  \node[annot, anchor=south west, font=\scriptsize]
    at (3.05,2.18) {4.2 reroot preview / 4.3 decoders};
  \node[annot, anchor=south west, font=\scriptsize, text=lv5]
    at (4.05,2.70) {4.3 rebuild-from-string};
  \node[annot, anchor=south west, font=\scriptsize, text=lv5]
    at (5.05,3.22) {scaled summit (834)};
\end{tikzpicture}
\caption*{Chapter~4 on the six-step ladder: the three topics carry you from raw storage
to rooting and decoding ($\lv4$), the rebuild-from-string topic reaches $\lv5$, and the
scaled summit (Sum of Distances, a reroot) sits on step~5. Beginner path $=$ steps 1--3.}
\end{figure}

% ---------------- The two running examples ----------------
\subsection*{The Two Running Examples}

A chapter about \emph{storage} needs \emph{two} concrete trees, because the whole point
is that a binary tree and a general tree want different arrays. Memorise both; every
figure, worked example and trace in this chapter uses them and nothing else.

\medskip
\noindent\textbf{Example $R$ -- the binary tree (a pointer/implicit-array case).}
This is the very six-node tree Chapters~2 and~3 recursed over and traversed; we now look
at \emph{how it sits in memory}. Its full definition and orders live in
\figref{c03-running}; here we only pin down its array views.

\begin{figure}[H]
\centering
\begin{tikzpicture}[diagram]
  \node[troot] (r) at (0,0) {$1$};
  \node[tnode] (a) at (-1.7,-1.4) {$2$};
  \node[tnode] (b) at (1.7,-1.4) {$3$};
  \node[tleaf] (c) at (-2.6,-2.8) {$4$};
  \node[tleaf] (d) at (-0.8,-2.8) {$5$};
  \node[tnode] (e) at (1.7,-2.8) {$6$};
  \draw[tedge] (r) -- (a); \draw[tedge] (r) -- (b);
  \draw[tedge] (a) -- (c); \draw[tedge] (a) -- (d);
  \draw[tedge] (b) -- (e);
  \node[annot, anchor=west, align=left, font=\scriptsize] at (3.2,-2.2)
    {$3$ has \emph{no left} child:\\implicit index $2\!\cdot\!1{+}1=3$ is a gap};
\end{tikzpicture}
\caption{Running example $R$ (binary). As pointers: each \code{TreeNode} holds
\code{left},\ \code{right}. If we instead \emph{number nodes by their array slot}
(root $=0$, left $=2i{+}1$, right $=2i{+}2$), the values land at slots
$1\!\to\!0,\ 2\!\to\!1,\ 3\!\to\!2,\ 4\!\to\!3,\ 5\!\to\!4,\ 6\!\to\!5$ \emph{only}
when the tree is complete; here node $3$'s missing left child makes slot~3 a
\emph{hole} -- the exact reason implicit indexing needs a complete tree.}\label{fig:c04-running-r}
\end{figure}

\medskip
\noindent\textbf{Example $G$ -- the general contest tree (an adjacency/parent-array case).}
Seven nodes, \emph{undirected} input edges, rooted at $0$ for discussion. This is the
shape Codeforces hands you, and every representation on the pointer side fails to hold
it compactly -- you need adjacency or a parent array.

\begin{figure}[H]
\centering
\begin{tikzpicture}[diagram]
  \node[troot] (n0) at (0,0) {$0$};
  \node[tnode] (n1) at (-3,-1.4) {$1$};
  \node[tnode] (n2) at (0,-1.4) {$2$};
  \node[tnode] (n3) at (3,-1.4) {$3$};
  \node[tleaf] (n4) at (-3.9,-2.8) {$4$};
  \node[tleaf] (n5) at (-2.1,-2.8) {$5$};
  \node[tleaf] (n6) at (3,-2.8) {$6$};
  \draw[tedge] (n0) -- (n1); \draw[tedge] (n0) -- (n2); \draw[tedge] (n0) -- (n3);
  \draw[tedge] (n1) -- (n4); \draw[tedge] (n1) -- (n5); \draw[tedge] (n3) -- (n6);
  \node[annot, anchor=west, align=left, font=\scriptsize] at (4.2,0.1)
    {undirected edges\\$\{0\text{--}1,0\text{--}2,0\text{--}3,$\\$1\text{--}4,1\text{--}5,3\text{--}6\}$};
  \node[annot, anchor=east, align=left, font=\scriptsize] at (-4.8,-2.6)
    {$\degop(0)=3,\ \degop(1)=3$\\(highest degree, but not\\the same as the root!)\\[2pt]
     leaves $\{2,4,5,6\}$, $\height(G)=2$};
\end{tikzpicture}
\caption{Running example $G$ (general tree, $7$ nodes, rooted at $0$). Its three
representations, all describing the \emph{same} object, appear in
\secref{c04-reps} and \secref{c04-rooting}: adjacency
\code{adj} $=\{\{1,2,3\},\{0,4,5\},\{0\},\{0,6\},\{1\},\{1\},\{3\}\}$;
parent array $\parent=[-1,0,0,0,1,1,3]$; subtree sizes $\size=[7,3,1,2,1,1,1]$.
The bracket encoding of this rooting, $0(1(4)(5))(2)(3(6))$, is where
\secref{c04-formats} begins.}\label{fig:c04-running-g}
\end{figure}

\medskip
\noindent Keep these three lines for $G$ in mind -- they are re-derived, never asserted,
in \secref{c04-rooting}:
\begin{center}
\renewcommand{\arraystretch}{1.2}
\begin{tabular}{@{}l p{0.66\linewidth}@{}}
\toprule
\textbf{View of $G$} & \textbf{Value (rooted at $0$)} \\
\midrule
$\parent[i]$   & $[-1,\,0,\,0,\,0,\,1,\,1,\,3]$ \\
$\size[i]$ (subtree node count) & $[7,\,3,\,1,\,2,\,1,\,1,\,1]$ \\
bracket encoding & $0(1(4)(5))(2)(3(6))$ \\
\bottomrule
\end{tabular}
\end{center}

% ---------------- Roadmap diagram ----------------
\begin{figure}[H]
\centering
\begin{tikzpicture}[diagram, node distance=5mm,
    top/.style={boxnode, fill=primaryblue!10, draw=primaryblue, text width=27mm,
                minimum height=10mm, font=\small},
    res/.style={annot, text width=27mm, anchor=north}]
  \node[top] (t1) {\textbf{4.1} Pointers vs Arrays};
  \node[top, right=of t1] (t2) {\textbf{4.2} Rooting};
  \node[top, right=of t2] (t3) {\textbf{4.3} Encodings};
  \foreach \a/\b in {t1/t2, t2/t3} \draw[flow] (\a) -- (\b);
  \node[res] at ([yshift=-2mm]t1.south) {four storage models; memory layout;
    parent arrays; N-ary children vectors};
  \node[res] at ([yshift=-2mm]t2.south) {orient edges with a parent guard;
    per-node rooted data; the reroot preview};
  \node[res] at ([yshift=-2mm]t3.south) {edge list / parent seq / bracket /
    level-order array; rebuild from a string};
\end{tikzpicture}
\caption*{The three topics, left to right. 4.1 asks \emph{where do the edges live};
4.2 asks \emph{which node is on top}; 4.3 asks \emph{what string did the judge hand me}.
They compose: you decode (4.3) into a structure (4.1), root it (4.2), then solve.}
\end{figure}

% ---------------- Topic table ----------------
\begin{center}
\renewcommand{\arraystretch}{1.25}
\setlength{\tabcolsep}{4pt}
\begin{tabular}{@{}p{0.25\linewidth} p{0.35\linewidth} c p{0.23\linewidth}@{}}
\toprule
\textbf{Topic} & \textbf{Key results} & \textbf{Levels} & \textbf{Solved problems} \\
\midrule
\secref{c04-reps}\ \ Pointers vs Arrays
  & four storage models; cache \& layout; parent arrays; N-ary children & $\lv1$--$\lv3$
  & \probref{c04-nary-depth}, \probref{c04-validate-array}, \probref{c04-inform} \\
\addlinespace
\secref{c04-rooting}\ \ Rooting an Unrooted Tree
  & parent-guarded DFS; $\parent/\depth/\size/\height$; reroot preview & $\lv3$--$\lv5$
  & \probref{c04-reorder-routes}, \probref{c04-subtree-label}, \probref{c04-min-height}, \probref{c04-highest-score} \\
\addlinespace
\secref{c04-formats}\ \ Encodings and Input Formats
  & edge list / parent seq / bracket / level-order-with-nulls; rebuild & $\lv2$--$\lv5$
  & \probref{c04-string-tree}, \probref{c04-recover-preorder} \\
\bottomrule
\end{tabular}
\end{center}

% ---------------- Problem recognition guide ----------------
\subsection*{Problem Recognition Guide --- When You See \emph{this}, Think \emph{that}}

Representation problems announce themselves by their \emph{input format} and by what
they ask you to \emph{navigate}. Match the phrasing; every row points at the section
that builds the tool.

\begin{center}
\renewcommand{\arraystretch}{1.3}
\setlength{\tabcolsep}{4pt}
\begin{tabular}{@{}p{0.40\linewidth} p{0.36\linewidth} l@{}}
\toprule
\textbf{When the problem says\ldots} & \textbf{\textcolor{deepblue}{\textbf{Think\ldots}}} & \textbf{Section} \\
\midrule
``$n$ nodes, node $i$ has children \code{leftChild[i]}, \code{rightChild[i]}'' / ``is this one valid tree?''
  & parallel child \emph{arrays}; check in-degree $=1$, one root, DFS counts all $n$ & \secref{c04-reps} \\
\addlinespace
``you are given \code{parent[]}, where \code{parent[i]} is the parent of $i$''
  & parent array $\Rightarrow$ build \code{children[]} lists, then BFS/DFS down & \secref{c04-reps} \\
\addlinespace
``an undirected tree with $n-1$ \emph{edges}, rooted at $0$ (or node $1$)''
  & adjacency list $+$ a parent-guarded DFS to orient it & \secref{c04-rooting} \\
\addlinespace
``for each node, count/aggregate over \emph{its subtree}''
  & root once, one DFS returns $\size$/$\text{agg}$ per node & \secref{c04-rooting} \\
\addlinespace
``the answer if the tree were rooted at \emph{every} node'' / ``reorder edges so all point to $0$''
  & rerooting: compute for one root, then $O(1)$ shift across each edge & \secref{c04-rooting} (\chref{c20}) \\
\addlinespace
``a string like \code{1(2)(3(4))}'' / ``preorder with \# marks'' / ``an array with nulls''
  & decode the \emph{encoding}; a stack or recursive splitter rebuilds pointers & \secref{c04-formats} \\
\bottomrule
\end{tabular}
\end{center}

\begin{center}
\renewcommand{\arraystretch}{1.25}
\begin{tabular}{@{}p{0.27\linewidth} p{0.45\linewidth} p{0.20\linewidth}@{}}
\toprule
\textbf{Constraint you see} & \textbf{What it forces} & \textbf{Chapter-4 tool} \\
\midrule
$N$ up to $2\cdot10^{5}$, general tree & adjacency as \code{vector<vector<int>>} or flat head/next arrays, not $\mathtt{map}$ & \secref{c04-reps} \\
tree may be a chain (depth $\approx n$) & recursion risks RE -- root it \emph{iteratively} or raise the stack & \secref{c04-rooting} \\
"answer for every root" in $O(n)$ total & naive re-root $O(n^2)$ -- take the reroot shift & \secref{c04-rooting} \\
input is one long string / one array with nulls & it is an \emph{encoding}, not a structure yet -- decode first & \secref{c04-formats} \\
\bottomrule
\end{tabular}
\end{center}

\begin{notebox}[Not this chapter?]
If the tree is \emph{already} a pointer \code{TreeNode} and the task is to \emph{visit}
it, the order is Chapter~3 (\chref{c03}), not representation. If a heap's
array-backed priority queue is the actual point, that is \chref{c07} -- this chapter only
teases the $i \leftrightarrow 2i{+}1$ layout. If you need \code{tin}/\code{tout} times
and \emph{subtree-as-range} for queries, that is the Euler tour (\chref{c24}); if you
need to \emph{answer from every root} at full generality, that is Rerooting (\chref{c20});
if the tree is defined by a \emph{Prufer code}, that is Counting Labelled Trees
(\chref{c32}). This chapter owns \emph{where the edges live and which node is the root}.
\end{notebox}

% ---------------- Notation ----------------
\subsection*{Notation Used in This Chapter}

\begin{center}
\renewcommand{\arraystretch}{1.2}
\begin{tabular}{@{}p{0.20\linewidth} p{0.27\linewidth}@{\quad}p{0.17\linewidth} p{0.27\linewidth}@{}}
\toprule
\textbf{Symbol} & \textbf{Meaning} & \textbf{Symbol} & \textbf{Meaning} \\
\midrule
$R,\ G$          & the running binary / general trees & $\code{adj}[v]$ & neighbour list of $v$ (undirected) \\
$\parent[v]$     & index of $v$'s parent ($-1$ at root) & $\children[v]$ & $v$'s children \emph{after rooting} \\
$\depth(v)$      & edges from the root down to $v$ & $\size[v]$  & node count of $\subtree(v)$ \\
$\height(T)$     & longest root-to-leaf path, in edges & $\degop(v)$ & undirected degree of $v$ \\
$\subtree(v)$    & the tree rooted at $v$ & $\lv{3}$ & difficulty badge (runs $\lv{1}$--$\lv{6}$) \\
$2i{+}1,\,2i{+}2$ & implicit heap-style child slots of node $i$ & $\mathit{par}[\,]$ & the parent array, indexed by node \\
\bottomrule
\end{tabular}
\end{center}

As in Chapters~1--3, $\height$ and $\depth$ are counted in \textbf{edges}
($\height(\nil)=-1$), general trees are stored as \code{vector<vector<int>> adj} with a
parent guard, every listing is complete $\cpp$ (\code{\#include <bits/stdc++.h>}, the
standard \code{TreeNode} where the tree is binary), and node numbering follows the
problem: LeetCode binary trees are \emph{unindexed} (pure pointers), contest general
trees are \emph{$0$- or $1$-indexed} by node id.
